\begin{textsample}{NEG Dim 1 – rn – Score 2.00 – rn/rn\_ef\_36.txt}  \label{ex:f1_neg_055}
INSTRUMENTOS E PROCEDIMENTOS DE AVALIAÇÃO - ANOS INICIAIS Avaliar se o estudante expressa, por meio da escrita, os conceitos trabalhados.
Por exemplo : desenhos ( 1º e 2º anos ) e relatórios ( 3º, 4º e 5º anos ). - Avaliar se o estudante se comportou de acordo com as diretrizes sugeridas durante as aulas ; respeitou os outros estudantes ; aceitou os resultados das atividades competitivas e individuais. - Criar, como instrumento de avaliação, um caderno de anotações do professor, em que seja relacionado se o aluno : - planeja ações que contribuem para a realização das atividades em sala ( 1º e 2º anos ) ; - experimenta os objetos de conhecimento valorizando cada um deles ( 1º e 2º anos ) ; - descreve, de forma oral e escrita, as atividades propostas ( 1º e 2º anos ) ; - planeja e utiliza estratégias para solução de problemas relacionados aos objetos de conhecimento ( 3º, 4º e 5º anos ) ; - realiza, com os colegas, a construção de objetos de conhecimento que favorecem a participação de todos nas aulas de Educação Física ( 3º, 4º e 5º anos ) ; - compreende e explica, de forma organizada, os objetos de conhecimento ( 3º, 4º e 5º anos ) ; - colabora para o desenvolvimento das atividades propostas ( 3º, 4º e 5º anos ).
INSTRUMENTOS E PROCEDIMENTOS DE AVALIAÇÃO - ANOS FINAIS Avaliar se o estudante expressa, por meio da escrita e da oralidade, os conceitos trabalhados.
Por exemplo : avaliação escrita, relatórios e apresentação de trabalhos em feira de ciências e na sala de aula. - Avaliar se o estudante se comportou de acordo com as diretrizes sugeridas durante as aulas ; respeitou os outros estudantes ; aceitou os resultados das atividades competitivas e individuais. - Criar, como instrumento de avaliação, um caderno de anotações do professor, em que seja relacionado se o aluno : - experimenta os objetos de conhecimento valorizando cada um deles ; - vivencia nas aulas de Educação Física os objetos de conhecimento ; - pratica, com ou sem materiais, atividades relacionadas aos objetos de conhecimento ; - constrói materiais recicláveis e implementos que colaborem para as atividades nas aulas de Educação Física ; - produz e cria regras que facilitem as atividades em sala ; - executa as atividades propostas pelo professor ; - identifica com clareza os objetos de conhecimento, suas semelhanças e as diferenças entre eles.

% matched lemmas: 
\end{textsample}
